\chapter{Archivos, pendrives y copias de seguridad}\label{cap:11-archivos}

\section{El gestor de archivos}

Tus archivos se manejan con el \termino{gestor de archivos}: en XFCE se
llama \emph{Thunar}; en otros escritorios, Caja (MATE), PCManFM-Qt
(LXQt), Archivos (GNOME) o Dolphin (KDE Plasma). Se abre desde el panel de
abajo, con el icono de la carpeta, o haciendo doble clic en \emph{Carpeta
personal} del escritorio.

\captura{thunar}{Thunar, el gestor de archivos de XFCE}

\subsection{Tu carpeta personal}
Cada usuario tiene su \emph{carpeta personal}, el equivalente a la carpeta
de usuario de Windows. Dentro están \emph{Documentos}, \emph{Descargas},
\emph{Imágenes}, \emph{Música}, \emph{Videos} y \emph{Escritorio}. Guarda
ahí lo tuyo: el resto del disco es del sistema y no hace falta tocarlo.

\subsection{Lo de todos los días}
\begin{itemize}
  \item \textbf{Copiar}: clic derecho sobre el archivo, \menu{Copiar}; ve a
        la otra carpeta, clic derecho en un espacio vacío, \menu{Pegar}.
        También se puede arrastrar.
  \item \textbf{Cambiar el nombre}: clic sobre el archivo y \keys{F2}.
  \item \textbf{Borrar}: \keys{Supr} lo manda a la \emph{Papelera}; desde
        ahí se puede recuperar hasta que la vacías.
  \item \textbf{Crear una carpeta}: clic derecho en un espacio vacío,
        \menu{Crear carpeta}.
  \item \textbf{Buscar}: \keys{\ctrl + F} y escribe parte del nombre.
  \item \textbf{Ver los archivos ocultos}: \keys{\ctrl + H}. Los archivos
        cuyo nombre empieza con un punto son configuraciones; mejor no
        tocarlos.
\end{itemize}

\subsection{Archivos comprimidos (.zip, .rar, .7z)}
Clic derecho sobre el archivo, \menu{Extraer aquí}. Para comprimir:
selecciona los archivos, clic derecho, \menu{Crear un archivador} (o
\menu{Comprimir}) y elige el formato \texttt{.zip}, el que se abre en
cualquier equipo.

\section{Pendrives, discos externos y el celular}

Al conectar un pendrive o un disco externo aparece solo en la columna de
la izquierda del gestor de archivos. Haz clic para abrirlo.

\begin{cuidado}
Antes de sacar el pendrive, haz clic en el triángulo de expulsar que
aparece al lado de su nombre (o clic derecho, \menu{Expulsar}) y espera a
que desaparezca. Si lo sacas mientras se está copiando algo, el archivo
y a veces todo el pendrive pueden quedar dañados.
\end{cuidado}

\begin{nota}
Copiar archivos grandes a un pendrive lento puede tardar: el sistema
copia primero a la memoria y después al pendrive. El Afinador de Lázaro
ajusta esto para que el equipo no se congele mientras tanto
(capítulo~\ref{cap:15-herramientas}). Espera siempre a expulsar.
\end{nota}

\subsection{El celular}
Conecta el celular por cable y, en Android, elige \emph{Transferencia de
archivos} en la notificación del USB. Aparece en el gestor de archivos y
puedes copiar las fotos de la carpeta \texttt{DCIM}. Con un iPhone,
desbloquéalo y toca \emph{Confiar} cuando lo pregunte.

\subsection{Formatear un pendrive}
Formatear borra todo y lo deja como nuevo. En el gestor de archivos, clic
derecho sobre el pendrive, \menu{Formatear} (si tu escritorio no lo
ofrece, instala \emph{Discos} desde la tienda). Elige \textbf{exFAT} si
lo vas a usar también en Windows, Mac o un televisor con archivos de más
de 4 GB; \textbf{FAT} para lo más compatible con equipos antiguos.

\section{Copias de seguridad}\label{sec:copias}

Un disco puede fallar sin avisar, sobre todo en un equipo viejo. La regla
que recomiendan los especialistas es simple: tener \textbf{al menos dos
copias} de lo importante, en lugares distintos \cite{incibe-copias}.

\begin{pasos}
  \item Consigue un disco externo o un pendrive grande que uses solo para
        esto.
  \item Una vez por semana (o cada vez que guardes algo importante),
        conéctalo y copia tus carpetas \emph{Documentos} e \emph{Imágenes}
        dentro de una carpeta con la fecha, por ejemplo
        \texttt{Copia 2026-11-09}.
  \item Expúlsalo y guárdalo en otro lugar.
\end{pasos}

\begin{consejo}
Si prefieres que se haga sola, la tienda de aplicaciones ofrece
\emph{Déjà Dup} (\enquote{Copias de seguridad}), que guarda una copia
programada en un disco externo o en la nube y permite recuperar versiones
anteriores de un archivo.
\end{consejo}

\fuentes
